\begingroup
\setlength{\tabcolsep}{4.5pt}
\renewcommand{\arraystretch}{1.28}
\begin{longtable}{@{}>{\raggedright\arraybackslash}p{.14\textwidth}>{\raggedright\arraybackslash}p{.33\textwidth}>{\raggedright\arraybackslash}p{.25\textwidth}>{\raggedright\arraybackslash}p{.22\textwidth}@{}}
\caption{Réalisation des composés, des secteurs topologiques et des états virtuels.}\label{tab:masas-realizacion-no-escalar}\\
\toprule
\textbf{Classe} & \textbf{Objeto HMT} & \textbf{Application de réalisation} & \textbf{Conclusion et conditions} \\
\midrule
\endfirsthead
\toprule
\textbf{Classe} & \textbf{Objeto HMT} & \textbf{Application de réalisation} & \textbf{Conclusion et conditions} \\
\midrule
\endhead
\midrule
\multicolumn{4}{r}{\footnotesize Suite à la page suivante.}\\
\endfoot
\bottomrule
\endlastfoot
mésons & Enregistrement double, frontière de liaison et signature composée. & Enregistrement composé et liaison $\to$ carte scalaire. & Composition des enregistrements ; la liaison précède la carte scalaire. \\
baryons & Enregistrement triple, retenue et signature composée. & Enregistrement composé et liaison $\to$ carte scalaire. & Composition des enregistrements ; la liaison précède la carte scalaire. \\
Fibonacci $(1,\tau)$ & Anneau de fusion de Fibonacci : $\tau\otimes\tau=1+\tau$. & Fusion et tressage $\to$ observable physique du secteur. & Codomaine exact de fusion et de tressage ; la comparaison utilise ses invariants propres. \\
Ising $(1,\psi,\sigma)$ ; réalisation de Majorana séparée & Anneau de fusion d’Ising : $(1,\psi,\sigma)$. & Fusion et tressage $\to$ observable physique du secteur. & Codomaine exact de fusion et de tressage ; la comparaison utilise ses invariants propres. \\
sectores $\mathbb Z/6$ & Caractère de tressage sur $\mathbb Z/6$. & Fusion et tressage $\to$ observable physique du secteur. & Codomaine exact de fusion et de tressage ; la comparaison utilise ses invariants propres. \\
secciones virtuales & Système inverse d’extensions et horizon de persistance. & Prolongement enregistré $\to$ horizon physique de persistance. & Objet exact de prolongement et d’obstruction ; sa sortie est l’horizon de persistance. \\
\end{longtable}
\endgroup
